% ============================================================
% VI. CONCLUSIONES
% ============================================================

\section{Conclusiones}

La Transformada Rápida de Fourier, calculada mediante \texttt{signal\_lab.spectral.compute\_fft}, permitió transformar las ocho señales de audio del dominio temporal al dominio frecuencial, facilitando la identificación de componentes espectrales característicos de cada muestra dentro de la ventana de análisis de 3 s.

En los sonidos animales, procesados a partir de grabaciones reales de dominio público/CC, se observaron diferencias claras y cuantificables entre las frecuencias dominantes del gato ($1545.17$ Hz), el perro ($459.73$ Hz) y el gallo ($1731.67$ Hz). En los instrumentos musicales, también reales, el contrabajo dio la frecuencia dominante más baja ($66.00$ Hz) y el piano la más alta ($1179.33$ Hz) —por encima incluso de la flauta ($546.33$ Hz)—, lo que ilustra que la frecuencia dominante depende tanto del instrumento como del contenido musical específico capturado en la ventana analizada (una nota vs. una frase). En las voces —aquí sintetizadas, no grabadas— también se observaron diferencias espectrales claras entre los dos motores de síntesis empleados.

Los siete pares de señales comparables (tres de animales, tres de instrumentos, uno de voces) resultaron, sin excepción, ``claramente separables'' según el puntaje de separabilidad de la Ec.~\ref{eq:separabilidad}. Esta uniformidad, lejos de ser una confirmación general de que la frecuencia dominante distingue bien cualquier par de señales, expuso una debilidad de diseño del propio puntaje (Sección~\ref{sec:discusion-unidades}): al dividir una diferencia de frecuencia (Hz) por una suma de desviaciones estándar de amplitud (adimensional, típicamente pequeña para audio normalizado), el denominador es sistemáticamente pequeño y el umbral de separabilidad se supera casi siempre, perdiendo poder discriminante justo cuando más se necesita —para distinguir separaciones grandes de separaciones marginales.

La media y la desviación estándar, calculadas de forma aislada, no demostraron ser características suficientes para identificar las diferentes señales; solo adquieren utilidad como parte del denominador de un puntaje de separabilidad, y ese puntaje, tal como está definido, requiere corrección antes de usarse como criterio cuantitativo definitivo.

Finalmente, los resultados de animales e instrumentos ya están respaldados por audio real, aunque con una sola grabación por categoría; los de voces siguen siendo exploratorios por depender de síntesis y no de grabaciones humanas. Los próximos pasos naturales son: (1) sustituir las voces sintetizadas por grabaciones humanas reales, siguiendo el mismo procedimiento documentado en \texttt{audio\_samples/SOURCES.md}; (2) recolectar varias grabaciones independientes por categoría; (3) corregir la Ec.~\ref{eq:separabilidad} para que ambos términos compartan unidades comparables (por ejemplo, normalizando por el ancho de banda de Nyquist); y (4) incorporar características espectrales adicionales (armónicos, energía por bandas) antes de evaluar un clasificador automático de fuentes sonoras.
